\documentclass[../main.tex]{subfiles}
\graphicspath{{\subfix{../images/}}}

\begin{document}

\section{Introduction and Context of Autonomous Maritime Perception}

\subsection{Evolution of Maritime Perception Systems}

Real-time perception for autonomous surface vessels has evolved substantially over the past decade, primarily driven by advances in deep learning architectures and sensor technologies.
Early maritime perception systems relied on rule-based processing of radar and sonar returns, suitable for detecting large vessels but inadequate for identifying diverse maritime objects or extracting semantic information.
The emergence of convolutional neural networks trained on large-scale terrestrial datasets demonstrated that learned representations could achieve superior performance compared to hand-engineered features, prompting the adaptation of these methods to maritime applications.

However, direct transfer of automotive perception methods to maritime environments proved problematic.
Automotive benchmarks such as KITTI~\cite{geiger2013} and Waymo~\cite{sun2020} assume dense \ac{LiDAR} ground plane returns, structured road environments, and relatively predictable lighting conditions---assumptions that do not hold for maritime operations.
Recent surveys of maritime autonomy~\cite{bae2023, zhang2021, xue2025, ferreira2022, bai2022} reinforce the perception of perception as a critical enabling technology, yet systematic performance characterization of detection methods under maritime operational conditions remains limited.

\subsection{Unique Challenges of the Maritime Environment}

The maritime environment presents distinct challenges compared to terrestrial scenarios, fundamentally affecting both vision-based and \ac{LiDAR}-based perception systems.
Studies of visual detection in coastal and harbor settings report substantial performance degradation due to water reflections, high dynamic range, and low-texture backgrounds~\cite{prasad2017}.
Maritime scenes frequently exhibit extreme illumination variations: intense glare on water surfaces exceeds the capture capabilities of standard cameras, while low-contrast vessel silhouettes against sky or water backgrounds reduce detection confidence.
Specular reflections from wave surfaces introduce dynamic false targets and occlude true maritime objects, degrading detection reliability.

Platform dynamics further compound these challenges. 
Unlike ground vehicles, surface vessels experience continuous wave-induced motion. 
These motions continually alter sensor orientation, requiring motion compensation to maintain consistent spatial alignment across frames.

\ac{LiDAR} systems face complementary difficulties in maritime environments.
Water surfaces produce specular reflection, yielding only 10\% return rates compared to near-100\% return rates on asphalt~\cite{kunz2005}, creating sparse point clouds that lack the dense ground plane context exploited by automotive algorithms.
Atmospheric scattering from salt spray and humidity reduces effective \ac{LiDAR} range, while wave-induced motion requires IMU-based motion compensation with associated computational overhead~\cite{roriz2022}.

This review synthesizes research across vision and \ac{LiDAR}-based detection methods independently before addressing multimodal fusion architectures and calibration requirements, and concludes by identifying specific research gaps motivating this dissertation's contributions.

\section{Modern Vision-Based Object Detection for Autonomy}

\subsection{Deep Learning Foundations and Single-Stage Detectors}

The application of convolutional neural networks to visual recognition gained widespread adoption following Krizhevsky et al.'s demonstration that deep architectures trained on large-scale datasets could achieve superior performance on ImageNet~\cite{krizhevsky2017}.
Subsequent architectural innovations, including residual learning frameworks~\cite{he2016}, demonstrated that network depth correlates with representational capacity for complex visual recognition tasks.
Early object detection methods adapted these classification networks into two-stage pipelines.
Girshick et al.~\cite{girshick2014} introduced R-CNN, which applies convolutional neural networks to region proposals.
However, the computational overhead limited throughput.
Ren et al.~\cite{ren2016} addressed these limitations with Faster \ac{R-CNN}, introducing region proposal networks that enabled end-to-end training and substantially improved inference speed.

Single-stage detectors emerged to address latency constraints by reformulating detection as direct regression from image features to object classes and bounding box coordinates.
The YOLO family exemplifies this paradigm, enabling real-time operation through unified spatial prediction~\cite{redmon2016}.
Subsequent YOLO iterations introduced architectural refinements, including multi-scale detection heads and anchor-free prediction mechanisms \cite{terven2023}.

YOLOv8, released in 2023 as a commercial implementation by Ultralytics \cite{ultralytics}, introduced architectural innovations addressing limitations of earlier single-stage detectors.
The anchor-free detection strategy employs dual detection heads that independently identify object presence and predict bounding box coordinates, simplifying transfer learning and enhancing adaptability to varying aspect ratios.
Modern loss functions, including Complete Intersection-over-Union (CIoU) and Distribution Focal Loss, improve localization precision and boundary regression stability~\cite{redmon2016}.

Max pooling and image down-sampling contribute to CNN computational efficiency but impose a fundamental constraint on \textbf{Small Object Detection (SOD)}: convolutional architectures employing progressive down-sampling through pooling layers reduce spatial resolution at deeper network layers, causing objects smaller than approximately 16-32 pixels to produce low-confidence or unstable detections~\cite{rekavandi2022}.
This constraint is particularly problematic for maritime perception, where detection targets such as navigation buoys or distant vessels may occupy minimal pixel area even at operationally relevant ranges.
For maritime applications processing imagery at $640 \times 640$ pixel inputs---the standard resolution for YOLOv8 deployment---down-sampling corresponds to smaller or distant targets near the horizon occupying less than 0.5-1\% of the image frame, placing them at or below the effective detection threshold imposed by network architecture.
These architectural properties make YOLOv8 suitable for real-time use; however, the model's constraints on SOD, along with other environmental factors, motivate the investigation of complementary sensing modalities.

\subsection{Maritime-Specific Detection Challenges}

Visual detection in maritime environments presents challenges that are fundamentally distinct from those in terrestrial scenarios.
Prasad et al.~\cite{prasad2017} identify background motion from waves and wakes, low contrast objects, and environmental degradation as primary failure modes.
High dynamic range presents persistent challenges for maritime detection.
While HDR imaging offers theoretical advantages for high-contrast maritime environments, empirical studies report counterintuitive results~\cite{landaeta, liebergall, wang2019a}.
Networks pre-trained on standard dynamic range (SDR) datasets achieve 10\% higher recall than HDR-based approaches~\cite{landaeta, liebergall}.
This performance gap stems from pre-training bias: networks learn feature extraction strategies optimized for SDR characteristics, and the scarcity of HDR training data prevents effective adaptation through transfer learning.
Atmospheric conditions impose additional constraints: Zhang et al.~\cite{zhang2022a} addressed foggy maritime detection through preprocessing and attention mechanisms. In contrast, Rekavandi et al.~\cite{rekavandi2022} note that objects occupying less than 10\% of image area yield insufficient information for reliable classification.
Maritime scenarios exacerbate this difficulty, as vessels and navigation markers often appear near the horizon, occupying a minimal pixel area.
Shao et al.~\cite{shao2022} addressed this through multi-scale detection strategies and deformable convolution modules.
Maritime detection research remains fragmented, lacking standardized large-scale datasets analogous to KITTI or COCO~\cite{zhang2021, yang2024, su2023}.

\section{LiDAR-Based Object Detection for Autonomous Systems}

\subsection{Deep Learning Approaches for Point Cloud Processing}

\acl{LiDAR} sensors provide three-dimensional geometric measurements independent of illumination by measuring the time-of-flight of laser pulses.
Processing point cloud data presents unique computational challenges: unlike images with a regular grid structure, point clouds present non-uniform three-dimensional sampling of the surrounding area~\cite{zhou2018, shi2019}.
This sparse sampling presents opportunities for a variety of methods of detection.

Deep learning has emerged as the dominant paradigm for point cloud object detection, with neural networks designed to learn hierarchical feature representations from three-dimensional data.
These learned approaches leverage large labeled datasets to extract spatial patterns and semantic information directly from point coordinates.
The following architectural families represent the evolution of neural network-based methods, progressing from point-based processing to voxelization strategies and hybrid architectures.

Point-based networks (PointNet~\cite{qi2017}) process raw coordinates directly; however, they scale unfavorably with point cloud size~\cite{shi2019}.
Voxelization methods address scalability by discretizing point clouds into regular grids.
Zhou and Tuzel~\cite{zhou2018} developed VoxelNet, which transforms sparse point clouds into dense volumetric representations.
Yan et al.~\cite{yan2018} extended this with SECOND, introducing sparse convolutional operations that improve efficiency by processing only occupied voxels.
Lang et al.~\cite{lang2019} proposed PointPillars, which encodes point clouds into vertical columns, enabling efficient 2D convolutional processing.
The fundamental trade-off in voxelization is quantization error: small objects or fine geometric details may be lost during the discretization process~\cite{zhou2018}.
Hybrid architectures (PV-RCNN~\cite{shi2019, shi2020}) combine the strengths of point-based and voxel-based processing, achieving state-of-the-art automotive accuracy but assuming dense ground plane returns that do not hold for maritime applications where water surfaces produce sparse returns~\cite{kunz2005}.

Recent innovations demonstrate continued advancement, but with substantial computational requirements.
BEVFusion~\cite{liang2022, liu2023b} achieves state-of-the-art accuracy but requires desktop-class GPUs with substantial power consumption and achieves only 15-20 frames per second~\cite{liu2023b}---this power budget alone exceeds total available power for many small autonomous vessels.
Transformer-based architectures that apply self-attention to point clouds offer theoretical advantages but require multi-GPU configurations with substantial power consumption and have not been validated under maritime-specific conditions~\cite{xie2024}.

Beyond computational constraints, these neural network-based approaches face limitations for safety-critical maritime applications.
Enforcing strict rules in neural networks, such as ensuring object separation distances, is difficult due to their non-deterministic nature.
This limitation is particularly problematic when using spatial data for navigational purposes and obstacle avoidance~\cite{coyleE}.
Given this persistent challenge, there is a continued need to investigate deterministic approaches that can provide stricter guarantees for safety-critical applications.

\subsection{Deterministic Geometric Alternatives}

In contrast to neural network-based approaches, deterministic geometric algorithms offer advantages for resource-constrained embedded platforms, where runtime predictability and minimal training data requirements are highly valued.

\textcolor{red}{(Need: Historical point cloud clustering methods - DBSCAN, Euclidean clustering, region growing, or similar precursor geometric segmentation algorithms that established the foundation for grid-based approaches)}

Building on these clustering foundations, Coyle~\cite{coyleE} developed Grid-Based Clustering and Concave Hull Extraction (GB-CACHE), a deterministic approach specifically designed for unmanned systems perceiving objects on flat surfaces, including water.
GB-CACHE efficiently segments point clouds through grid-based spatial partitioning, then extracts concave hull boundaries that accurately represent objects with irregular geometries.
Computational efficiency analysis demonstrates real-time feasibility on low to mid-grade computing platforms~\cite{coyleE}.

The deterministic nature of GB-CACHE provides specific operational advantages.
Processing latency remains bounded and predictable, eliminating variable inference times characteristic of neural networks, where batch size and scene complexity affect throughput.
Rule enforcement for spatial constraints---critical for obstacle avoidance and navigation---can be rigorously guaranteed, whereas neural approaches provide probabilistic outputs without strict guarantees~\cite{coyleE}.
Training data requirements are minimal; the algorithm requires only geometric parameters rather than extensive labeled datasets---a practical advantage for maritime deployment where training data scarcity persists.

\subsection{Maritime-Specific LiDAR Challenges}

Maritime \ac{LiDAR} performance differs fundamentally from automotive scenarios.
Automotive systems achieve near 100\% return rates on asphalt, providing dense point clouds with abundant ground plane context~\cite{roriz2022}.
Water surfaces produce specular reflection, yielding only 10\% return rates versus near-100\% on asphalt~\cite{kunz2005, roriz2022}, creating sparse point clouds lacking dense ground plane context exploited by automotive algorithms~\cite{zhou2018, shi2019}.
Kunz et al.~\cite{kunz2005} measured return rates over calm water as low as 10\% of the emitted pulses, noting that water surface characteristics produce significantly weaker returns than small floating targets, creating a favorable contrast for target detection despite a sparse environmental context.

The sparsity of water returns poses challenges to detection algorithms designed for dense point clouds in automotive applications.
Maritime applications lack dense reference, necessitating alternative spatial reasoning strategies when objects appear isolated without contextual grounding.
Atmospheric effects reduce effective range by 50-70\% in heavy haze~\cite{roriz2022}, while wave-induced motion requires IMU-based motion compensation, introducing processing latency~\cite{xie2024, ahmed2024}.
Platform motion caused by continuous wave-induced pitch, roll, and heave motions vary dynamically with sea state.
Ahmed et al.~\cite{ahmed2024} developed specialized Kalman filtering approaches tailored for maritime platform motion, reporting 25-30\% improvement in detection accuracy when compensating for tilt-induced point cloud distortions.

Recent maritime \ac{LiDAR} detection work demonstrates growing interest yet reveals substantial validation gaps.
Xie et al.~\cite{xie2024} achieved an overall detection accuracy of 74.1\% by integrating state-of-the-art networks for ship detection, although performance degraded for vessels under 5 meters and at ranges exceeding 40 meters.
Detection methods designed for dense ground plane returns require fundamental adaptation for sparse maritime point clouds where objects lack grounding context~\cite{kunz2005}.
A systematic evaluation of detection performance versus water surface characteristics, motion compensation quality, and embedded platform profiling (including latency, power, and throughput) is absent from the reviewed literature~\cite{xie2024, ahmed2024}.

The complementary limitations of camera and \ac{LiDAR} modalities---vision failing under extreme lighting while LiDAR struggles with sparse returns---suggest potential benefits from sensor fusion.

\section{Sensor Fusion Paradigms}

Sensor fusion combines information from multiple modalities to overcome the limitations of individual sensors.
Recent surveys~\cite{feng2021, tufekci2023, huang2024} identify fusion level---early, mid, or late---as the primary architectural distinction determining computational requirements, calibration sensitivity, and failure mode propagation relevant to maritime deployment.
Each fusion paradigm offers distinct trade-offs between peak accuracy potential, robustness to sensor degradation, and real-time implementation feasibility~\cite{feng2021}.

\subsection{Early Fusion: Data-Level Integration}

Early fusion combines raw sensor data before feature extraction, projecting measurements from different modalities into a common representation~\cite{cui2022}.
Cui et al.~\cite{cui2022} note that early fusion enables networks to learn joint features during training.
PointPainting exemplifies this paradigm by projecting \ac{LiDAR} points into camera image coordinates, assigning semantic labels from segmentation networks, then processing the semantically enriched point cloud through object detectors~\cite{cui2022}.

However, early fusion imposes stringent calibration and synchronization requirements, creating practical deployment challenges.
Projection errors of even a few degrees of rotation or centimeters of translation misalign features, corrupting the joint representation~\cite{cui2022}.
Temporal synchronization must be precise; misaligned timestamps create spatial inconsistencies, degrading fusion quality.
Sensor-specific noise and failures propagate directly: corrupted camera images from lens occlusion or extreme lighting produce degraded semantic labels affecting all projected \ac{LiDAR} points~\cite{cui2022}.
For maritime platforms experiencing continuous wave-induced motion and intermittent sun glare, calibration sensitivity and failure propagation of early fusion represent significant operational risks.
Maritime platform dynamics introduce additional calibration complexity beyond automotive scenarios.
While automotive research established general calibration tolerances~\cite{cui2022}, wave-induced motion can create time-varying calibration errors even when initial alignment is precise.
These dynamic calibration challenges affect all fusion paradigms but with varying severity depending on the required precision of spatial correspondence.

\subsection{Mid Fusion: Feature-Level Integration}

Mid fusion extracts features independently from each modality before combining them in intermediate network layers~\cite{huang2024}.
Huang et al.~\cite{huang2024} systematically compare fusion paradigms, finding that mid fusion balances accuracy and robustness: independent feature extraction prevents raw sensor corruption from propagating between modalities, yet learned fusion layers capture cross-modal correlations unavailable to late fusion.

BEVFusion~\cite{liang2022, liu2023b} transforms camera features into bird's-eye-view representations compatible with \ac{LiDAR} BEV features, combining them through efficient convolution operations.
Liu et al.~\cite{liu2023b} note that BEVFusion achieves state-of-the-art performance but requires desktop-class GPU hardware, exceeding embedded platform capabilities by an order of magnitude.
For battery-powered maritime platforms with constrained power budgets, mid-fusion's computational demands preclude direct deployment without substantial architectural optimization.

\subsection{Late Fusion: Decision-Level Integration}

Late fusion operates on independent object detections from each modality, combining outputs through spatial association~\cite{wang2020a, pang2020}.
This approach maintains modular detection pipelines, fusing results at the decision level after individual sensor processing is complete.
While sacrificing joint feature learning enabled by earlier fusion, late fusion offers architectural properties advantageous for resource-constrained maritime deployments.

Wang et al.~\cite{wang2020a} survey multi-sensor fusion in automated driving, emphasizing that late fusion enables graceful degradation: sensor failures affect only individual detection streams without corrupting fusion outputs.
Pang et al.~\cite{pang2020} developed CLOCs, which fuse detection candidates before non-maximum suppression, training networks to predict fused detection quality based on geometric and semantic consistency between modalities.

Late fusion's modular architecture provides several practical advantages for maritime deployment.
Computational requirements scale approximately linearly with sensor count rather than exponentially as with joint processing~\cite{wang2020a}.
Processing latency remains predictable because each modality follows fixed pipelines without runtime cross-modal dependencies during feature extraction.
Adding or removing sensors requires minimal system redesign---new detection streams integrate through updated fusion modules without modifying existing pipelines~\cite{feng2021}.
Most critically for maritime applications, individual sensor failures degrade performance gracefully rather than causing catastrophic fusion failure~\cite{huang2024}.
Temporal synchronization requirements are relaxed compared to early fusion: late fusion spatial association operates on object-level regions rather than pixel-level correspondences, tolerating larger timing offsets---typically 50 to 100 milliseconds versus 10 milliseconds for early fusion~\cite{huang2024}.
For maritime platforms where precise hardware synchronization may be impractical and wave-induced motion creates apparent object movement, relaxed timing tolerances represent operational advantages.

\subsection{Comparative Analysis and Maritime Fusion Gaps}

Huang et al.~\cite{huang2024} systematically evaluate fusion paradigms under simulated sensor degradation and synchronization error.
Late fusion maintains superior robustness when individual sensors fail or produce degraded outputs through graceful degradation.
However, under ideal conditions with clean data and perfect calibration, mid fusion achieves higher peak accuracy by learning cross-modal correlations unavailable to late fusion~\cite{huang2024}.
Wang et al.~\cite{wang2020a} note that mid-fusion architectures approximately double processing requirements compared to single-modality detection.
For battery-powered maritime platforms, doubled power consumption directly reduces mission endurance.

Table~\ref{tab:fusion_comparison} summarizes the key trade-offs between fusion paradigms relevant to maritime USV deployment under power and computational constraints.

\begin{table}[htbp]
\centering
\caption{Comparison of Fusion Paradigms for Maritime USV Deployment}
\label{tab:fusion_comparison}
\small
\begin{tabular}{lccccc}
\hline
\textbf{Paradigm} & \textbf{Computational} & \textbf{Calibration} & \textbf{Timing} & \textbf{Sensor} & \textbf{Peak} \\
                  & \textbf{Demand} & \textbf{Sensitivity} & \textbf{Tolerance} & \textbf{Failure} & \textbf{Accuracy} \\
\hline
Early Fusion & High & Very High & $\sim$ 10 ms & Propagates & High \\
             &      & ($<1^\circ$ rotation) & (tight) & globally &  \\
Mid Fusion & Very High & High & $\sim$20ms & Partial & Very High \\
           & ($\sim$2$\times$ baseline) & ($<1^\circ$ rotation) & (moderate) & propagation &  \\
Late Fusion & Low & Low & 50-100ms & Graceful & Moderate \\
            & ($\sim$1.1$\times$ baseline) & ($>5^\circ$ tolerant) & (relaxed) & degradation &  \\
\hline
\end{tabular}
\end{table}

These trade-offs inform the selection of late fusion as the primary focus of this dissertation, prioritizing robustness and computational efficiency over peak accuracy under ideal conditions.
For small autonomous vessels operating under power constraints and experiencing continuous platform motion, late fusion's tolerance to calibration errors, relaxed synchronization requirements, and graceful degradation characteristics align with operational requirements~\cite{huang2024, wang2020a}.

Automotive fusion methods~\cite{feng2021, cui2022, huang2024} evaluate performance on terrestrial datasets where environmental characteristics differ fundamentally from maritime scenarios---water reflections, extreme dynamic range, sparse \ac{LiDAR} returns, and platform motion require empirical validation.
Computational requirements, power consumption, and latency characteristics on representative embedded hardware are rarely reported~\cite{liu2023a}.
Systematic evaluation sweeping timing offsets to establish acceptable synchronization tolerances for maritime platforms is absent~\cite{huang2024}, and wave-induced motion creates time-varying calibration errors differing from static errors typical of ground vehicles~\cite{cui2022}.
Most fusion studies report performance without rigorous comparison against optimized single-modality baselines~\cite{farahnakian2020, haghbayan2018, clunie2021}.

\section{Enabling Technologies: Calibration and Synchronization}

Accurate fusion depends on establishing spatial and temporal correspondence between sensors.
Spatial calibration defines rigid-body transformation relating sensor coordinate frames, while temporal synchronization ensures observations represent the same moment in time.
Precision requirements vary with the fusion paradigm---early fusion demands tighter tolerances than late fusion due to the need for pixel-level data alignment.

\subsection{Spatial Calibration Requirements}

\subsubsection{Intrinsic Calibration}

Camera intrinsic parameters---including focal length, principal point, and lens distortion coefficients---must be accurately estimated to enable geometric correspondence between sensors.
The pinhole camera model provides the mathematical foundation for projecting three-dimensional points onto two-dimensional image planes.
Calibration typically employs structured targets such as checkerboards or AprilTags, where known geometry enables parameter estimation through least-squares optimization.
\textcolor{red}{(Need: Zhang's camera calibration method - seminal work on planar calibration pattern approach for intrinsic parameter estimation)}
Multiple views of planar calibration patterns enable recovery of intrinsic parameters and distortion coefficients through least-squares optimization.
For maritime applications, lens distortion correction is particularly critical: wide-angle lenses commonly used for situational awareness introduce significant radial distortion that degrades detection accuracy if uncorrected.

\subsubsection{Extrinsic Calibration}

Extrinsic calibration establishes the geometric relationship between sensors, typically represented as a rotation matrix and translation vector defining the rigid-body transformation from one sensor's coordinate frame to another.
Traditional calibration employs fiducial targets visible to both sensors, solving for the transformation through least-squares optimization of corresponding point observations.
Manual calibration remains standard for many research systems, as it provides the precision necessary for reliable geometric correspondence.

Automated calibration methods have emerged as alternatives. Iyer et al.~\cite{iyer2018} developed CalibNet, which uses deep learning to predict calibration parameters directly from sensor data.
Xiao et al.~\cite{xiao2024} proposed CalibFormer, a transformer-based automatic LiDAR-camera calibration network~\cite{schneider2017, wu2021}.
However, these automated approaches require sufficient scene structure and may fail in textureless maritime environments, which are often dominated by sky and water.

Calibration accuracy directly affects fusion quality.
Rotation errors exceeding one degree or translation errors exceeding 10 millimeters significantly degrade early and mid-fusion performance.
These tolerances become more stringent as sensor separation increases.
For late fusion approaches operating on object-level detections rather than pixel-level correspondences, calibration tolerances are substantially relaxed: errors of several degrees or tens of centimeters may be acceptable depending on detection range and object size~\cite{huang2024}.

\subsection{Temporal Synchronization Protocols}

Temporal synchronization ensures observations from different sensors represent the same world state~\cite{westenberger2011}.
Timing offsets create spatial misalignment: a vessel moving at 5 meters per second experiences 0.5-meter displacement in 100 milliseconds, introducing apparent object motion that degrades fusion quality.
The required synchronization precision depends on platform velocity, sensor modality, and fusion approach.

\subsubsection{Network-Based Synchronization Standards}

Network Time Protocol (NTP) provides millisecond-level synchronization over standard Ethernet networks, sufficient for many late fusion applications where object-level association tolerates moderate timing errors~\cite{furgale2013}.
IEEE 1588 Precision Time Protocol (PTP) achieves sub-microsecond synchronization accuracy, enabling tight temporal alignment required for early and mid-fusion approaches~\cite{ptp2008}.
PTP employs hardware timestamping at network interface layers, eliminating variable software latencies that degrade NTP accuracy.

Westenberger et al.~\cite{westenberger2011} demonstrated that sensor-specific clock drifts and lost measurements (frame drops) introduce systematic timing errors that standard synchronization protocols cannot fully address.
Their approach detects frame drops and compensates for clock drift through continuous timestamp monitoring, achieving millisecond-level accuracy for automotive multi-sensor fusion.

For maritime platforms, wave-induced motion introduces an additional temporal challenge beyond static synchronization: IMU-based motion compensation requires precise timestamps to retrospectively correct sensor observations for platform orientation changes during data acquisition~\cite{liu2021}.

\subsubsection{High-Bandwidth Video Synchronization}

High-bandwidth video streams present unique synchronization challenges distinct from discrete sensor measurements.
Standard network-based approaches timestamp video frames at reception, introducing variable latencies from encoding, transmission, and buffering that corrupt temporal correspondence with other sensors.

Direct frame-level timestamp embedding within the video bitstream ensures temporal integrity throughout the acquisition and transmission pipeline.
The H.264/AVC video coding standard~\cite{wiegand2003} defines Supplemental Enhancement Information (SEI) messages as a mechanism for embedding metadata within the compressed video stream.
\textcolor{red}{(Need: Specific references on H.264/H.265 SEI NAL unit timestamp embedding methods, particularly Type 1 Picture Timing SEI messages or similar mechanisms for frame-level temporal annotation)}

For perception systems streaming high-resolution camera data over Ethernet to centralized processing units, SEI-based timestamp embedding addresses a critical gap: it preserves acquisition timestamps through the encoding, transmission, and decoding pipeline, enabling precise temporal alignment with \ac{LiDAR} point clouds for late fusion.
This methodology represents a practical engineering solution motivated by the operational requirement for temporal correspondence under network transmission constraints, as detailed in Chapter~\ref{realtime_object_detection}.

\section{Maritime Perception Datasets and Benchmarks}

Dataset availability fundamentally shapes algorithm development and validation.
Automotive perception benefits from large-scale benchmarks (KITTI~\cite{geiger2013}, Waymo, nuScenes~\cite{caesar2020}) providing millions of labeled examples with synchronized camera-LiDAR data.
These benchmarks enable quantitative comparison across algorithms and have accelerated progress in terrestrial autonomous vehicle perception.

Maritime perception lacks comparable resources; most datasets provide only monocular RGB imagery without depth information~\cite{su2023}.
Those including \ac{LiDAR} often lack precise temporal synchronization or provide only sparse annotations covering limited object categories.
Weather diversity remains underrepresented---most datasets capture clear conditions with few examples of rain, fog, or heavy seas.
Su et al.~\cite{su2023} provide a comprehensive survey of maritime vision datasets, documenting available resources and identifying critical gaps.
Zhang et al.~\cite{zhang2021} review perception technologies for autonomous ships, noting the scarcity of publicly available multimodal maritime datasets as a fundamental barrier to algorithm development.
Kim et al.~\cite{kim2022} improved the Singapore Maritime Dataset, providing RGB, infrared, \ac{LiDAR}, and radar modalities, though sparse ground truth limits the scope for training deep networks.
Recent maritime datasets~\cite{thompson2023} provide multimodal observations but remain limited in scale, weather diversity, and public availability.
Small datasets restrict model complexity without overfitting; limited weather coverage prevents robust evaluation; sparse annotations leave uncertainty about detection reliability~\cite{su2023}.

Comparative analysis with Unmanned Aerial System (UAS) perception research reveals relevant parallels for small USV platforms.
UAS datasets and fusion architectures address similar constraints---limited payload capacity, restricted power budgets, and real-time processing requirements on embedded hardware~\cite{thompson2023, haghbayan2018}.
However, UAS perception benefits from stable platforms and downward-looking sensing geometries, whereas USVs must contend with continuous wave-induced motion and forward-looking detection of low-contrast maritime objects against water backgrounds.
Despite these differences, UAS research on lightweight sensor fusion and embedded perception provides valuable methodological insights for maritime applications, particularly regarding power-constrained algorithm selection and real-time implementation strategies.

\section{Research Gaps and Dissertation Motivation}

Despite substantial progress in multimodal perception for autonomous systems, several critical gaps remain for maritime applications.
These gaps span algorithmic approaches, computational implementation, and empirical validation resources.

Most published fusion research emphasizes early or mid-level integration, with limited evaluation of decision-level approaches under maritime conditions.
Comparative studies in automotive contexts demonstrate that late fusion offers superior robustness under sensor dropout or environmental degradation---conditions frequently encountered during maritime operations~\cite{huang2024, wang2020a}.
However, these findings primarily derive from terrestrial datasets where environmental characteristics and sensor failure modes differ substantially from those in maritime scenarios.
Whether late fusion maintains similar robustness advantages under maritime-specific challenges, including water reflections, HDR lighting extremes, and platform motion, remains an open empirical question in the reviewed literature.
This dissertation addresses this gap through systematic evaluation of late fusion performance under maritime conditions, as described in Chapter~\ref{realtime_object_detection}.

% A critical gap exists in the establishment of \textbf{computational efficiency baselines} for maritime perception.
% No published work provides systematic latency and power consumption profiles for both single-modality and fusion approaches on representative embedded USV hardware under maritime operational conditions.
% This gap prevents informed algorithm selection for power-constrained platforms where computing loads directly affect mission endurance: without quantitative characterization of accuracy-latency-power trade-offs, practitioners cannot determine whether fusion provides sufficient accuracy improvement to justify increased computational demands.

Early and mid-fusion methods often exceed the computational budgets of embedded platforms, requiring desktop-class GPUs with typical power consumption of 200-400 watts for real-time operation---specifications incompatible with small, battery-powered autonomous vessels.
Late fusion's reduced computational requirements suggest potential feasibility for embedded deployment; however, systematic evaluation of accuracy-latency-power trade-offs on desktop-class maritime hardware is lacking in prior maritime fusion research.
This dissertation addresses these interconnected gaps through hardware profiling on Minion B workstation hardware (NVIDIA GeForce GTX 1080 GPU), establishing baseline performance for YOLO model variants (nano, small, medium) and GB-CACHE detection latency, power consumption, and accuracy characteristics.
Chapter~\ref{realtime_object_detection} details this systematic evaluation, documenting computational efficiency baselines that inform algorithm selection for power-constrained maritime platforms.

Few labeled multimodal maritime datasets exist with synchronized \ac{LiDAR} and camera observations under diverse conditions.
Existing datasets are either single-modality, which limits fusion research, lack precise synchronization, preventing temporal analysis, or provide limited weather and lighting diversity, thereby restricting generalization assessment.
This data gap directly impedes algorithm development and validation.
For this research, a synchronized \ac{LiDAR}-camera dataset was collected aboard the \acl{WAMV} platform during operational deployments, providing 778 annotated images with 1934 labeled maritime objects across diverse weather and lighting conditions for experimental validation, as described in Section~\ref{sec:sensor_data_dataset}.

This dissertation develops and validates a late-fusion framework for real-time \ac{LiDAR}-camera perception on embedded maritime platforms.
The approach is motivated by the advantages of late fusion for maritime deployment, including a modular architecture that isolates sensor failures, computational efficiency compatible with embedded hardware constraints, and robustness to calibration and synchronization errors compared to tightly coupled fusion approaches.
Research methodology quantitatively evaluates accuracy, timing, and computational efficiency through controlled experiments and real-world validation aboard the \ac{WAMV} research platform.

The following chapter describes experimental methodology employed to address these research questions, including detailed hardware configuration, sensor calibration procedures, data collection protocols, and late fusion algorithm design.

\end{document}
